\chapter{Sostenibilidad}
\label{chap:sostenibilidad}

\begingroup
\normalsize
\setstretch{1.05}
\setlength{\parskip}{0.3\baselineskip}

Este capítulo corresponde al hito inicial del informe de sostenibilidad. Se
completará en el hito final con el análisis ético y los consumos medidos.

\section{Autoevaluación}

Tras realizar la encuesta del proyecto EDINSOST, he podido identificar mis
puntos fuertes y débiles en las competencias medidas.

Mi punto más fuerte es el uso eficiente de los recursos. Este TFG trata
precisamente el compromiso entre precisión y coste computacional, y durante el
grado he aprendido a razonar sobre complejidad, memoria y tiempo de ejecución.
Sé reutilizar modelos preentrenados y datos públicos, y soy consciente del
coste energético de entrenar modelos grandes. También sé situar un problema
técnico en su contexto social: en este caso, la dificultad de registrar la dieta y su
relación con la salud pública.

Mis puntos débiles están en la cuantificación y en la práctica. Identifico de
forma cualitativa los impactos ambientales y sociales, pero no domino métodos
para medirlos, como la huella de carbono del entrenamiento, porque no he
participado en proyectos reales que lo requirieran. Del mismo modo, reconozco
los riesgos éticos de un sistema de estimación nutricional (privacidad, sesgo
de los datos y uso como consejo médico), pero no los he gestionado con
usuarios reales, ya que este proyecto investiga modelos y no crea un producto.

En este TFG mediré el cómputo empleado para estimar el impacto ambiental del
entrenamiento y documentaré los límites y riesgos del modelo, para evitar que
se use como herramienta clínica. Así trabajaré mis puntos débiles en
sostenibilidad.

\section{Dimensión económica}

El presupuesto del proyecto (capítulo~\ref{chap:gestion-economica}) asciende a
13\,008,52\,€, y el 73\,\% corresponde a mis horas de trabajo. Considero que es
un coste adecuado, ya que los recursos materiales se limitan a mi portátil,
software gratuito, cómputo alquilado solo cuando hace falta y un espacio de
trabajo. Si hubiera que reducirlo, el ahorro vendría sobre todo de las horas:
una búsqueda de hiperparámetros más acotada o probar menos modelos reduciría
tanto mi tiempo como el gasto en cómputo.

Al revisar el estado del arte, he visto que las aplicaciones que estiman
nutrientes a partir de fotografías no publican su coste de ejecución, aunque su
arquitectura lo condiciona. Cal AI usa modelos de Anthropic y OpenAI
\citep{PhotoCalorieApp}, por lo que cada fotografía implica una
inferencia en servidores externos. MyFitnessPal, en cambio, se basa en el kit de
Passio \citep{MyFitnessPalNutritionAI}, que puede ejecutar los modelos en el
propio dispositivo \citep{PassioConfigurationClassNutrition_ai} o en un servidor. Creo
que ejecutar un modelo ligero en el móvil es la opción más sostenible
económicamente: elimina el coste por predicción y el de mantener un servidor
de inferencia, que es elevado con los modelos grandes accesibles por API.
Además, el coste deja de crecer con el número de usuarios. A cambio, el modelo
puede perder precisión, y es justamente lo que este trabajo quiere cuantificar.

\section{Dimensión ambiental}

Considero que la realización de este proyecto tiene un impacto ambiental
reducido. El consumo de mi portátil está acotado a un máximo de 35,1\,kWh
durante las 540 horas, y el impacto principal será el de las GPU en la nube,
que calcularé al final a partir de las horas de GPU consumidas y de la potencia
de la GPU utilizada. Para minimizarlo reutilizo recursos que ya existen: un
conjunto de datos público (Nutrition5k), pesos preentrenados en ImageNet en
lugar de entrenar desde cero y el portátil que ya tenía. Además, pruebo cada
experimento con una pequeña parte de los datos antes de lanzarlo completo, no
repito ejecuciones ya terminadas y acoto la búsqueda de configuraciones.

Indirectamente, el proyecto también podría beneficiar al medio ambiente. Un
modelo ligero ejecutado en el dispositivo no necesita enviar cada imagen a un
servidor, por lo que espero que consuma bastante menos energía que la
inferencia en un centro de datos, sobre todo frente a modelos grandes. No
mediré el consumo de esas alternativas, pero sí el tamaño, la latencia y la
memoria de los modelos, que dan una idea de su coste energético.

\section{Dimensión social}

En el plano personal, este TFG me permite aprender a diseñar un protocolo
experimental riguroso, a entrenar y comparar modelos de visión entendiendo sus
arquitecturas y a desplegarlos en entornos con recursos limitados. También me
permite planificar y documentar un proyecto de investigación de principio a
fin, algo que considero muy valioso para mi futuro profesional.

En cuanto a la sociedad, la prevalencia de la obesidad en adultos se ha más que
duplicado desde 1990 \citep{ObesityOverweight}, y registrar la dieta a mano es
costoso y propenso a omisiones y errores de porción
\citep{whittonSystematicReviewExamining2022}. Estimar los nutrientes a partir
de una fotografía podría facilitar ese seguimiento y, si la inferencia es
local, evitaría enviar imágenes de las comidas a terceros. Considero que existe
una necesidad real: hay interés por la alimentación saludable
y aplicaciones comerciales que ofrecen esta
función \citep{HowUseMeal2026,Healthify,CalAIFood}, pero su exactitud no es
pública. Una comparación abierta como la de este trabajo aporta evidencia sobre
el compromiso entre error y coste.

\endgroup
